\begin{frame}{FAQ: 완화인가, 해결인가, GRU는 못나나, 셀 vs 은닉}
  \begin{block}{Q. LSTM/GRU를 쓰면 그래디언트 소실이 완전히
  사라지나요?}
    \textbf{아니다} --- ``\textbf{완화}이지 ``해결''이 아니다.
    게이트가 곱셈 경로 옆에 \textbf{덧셈 경로}를 추가해서 소실
    속도를 크게 늦추지만, \textbf{수천 토큰 이상}이면 LSTM/GRU도
    여전히 먼 과거 정보를 점점 잃어감.
    \kq{Week 12 Attention이 ``순차적으로 압축해서 기억한다''는
    접근 자체를 버리고 ``필요할 때마다 과거 전체를 직접 다시
    들여다본다''로 전환하는 이유가 바로 이것.}
  \end{block}
  \vskip0.5em
  \begin{block}{Q. GRU는 LSTM보다 항상 나쁜가요(게이트가 더 적으니까)?}
    \textbf{아니다} --- 파라미터가 적어(게이트 3개 대신 2개)
    \textbf{학습이 빠르고 데이터가 적을 때 과적합 위험도 낮음}.
    성능은 문제에 따라 둘 중 어느 쪽이 나을지 갈림 --- ``게이트가
    많을수록 무조건 좋다''는 직感は 성립하지 않음 (Week 6
    편향-분산 트레이드오프. 위 ``GRU 0.99 vs LSTM 0.86''이 직접
    증거).
  \end{block}
\end{frame}
